% Compléments techniques des chapitres de conception, d'architecture, de qualité, d'analyse et du
% tableau de bord.
%
% CE FICHIER EST UNE ANNEXE RÉDIGÉE, ET IL EST LU PAR LE CONTRÔLE DE PROVENANCE au même titre qu'un
% fichier de chapitre : `tests/test_provenance_des_chapitres.py` le nomme dans sa liste des fichiers
% d'annexe examinés, confrontée aux inclusions de `report/annexes.tex`. Le corps du rapport garde le
% raisonnement et les résultats ; le détail technique vit ici, avec ses étiquettes, et le corps y
% renvoie par `\ref`. Les sections suivent l'ordre des chapitres d'où vient leur matière.
%
% Ce fichier repose sur :
%   DOC: s22, s27, s30
%   OBS: REL-PAT.A01, REL-PAT.A03, REL-PAT.D01, REL-PAT.D02, REL-PAT.D03, REL-PAT.D05,
%   OBS: REL-PAT.D06, REL-PAT.D09, REL-PAT.D10, REL-PAT.D11, REL-PAT.H05, REL-PAT.H07,
%   OBS: REL-PAT.H12, REL-RDV.C01, REL-RDV.C02, REL-RDV.C05, REL-RDV.C07, REL-RDV.C08,
%   OBS: REL-RDV.R01, REL-RDV.R02, REL-RDV.R03, REL-RDV.R08, REL-RDV.R10, REL-RDV.R11,
%   OBS: REL-RDV.R12
%   HYP: (aucun)
%   CHI: artefact-part, blocage-nom-adresse, blocage-nom-telephone, blocage-parents-naissance, blocage-piece-identite, blocage-reduction-nom-adresse
%   CHI: blocage-reduction-nom-telephone, blocage-reduction-parents-naissance, blocage-reduction-piece-identite, blocage-reduction-union, colonnes-aval-decrites, colonnes-aval-total
%   CHI: colonnes-intermediate, colonnes-marts, colonnes-source-catalogue, colonnes-source-catalogue-doc, colonnes-source-catalogue-hyp, colonnes-source-catalogue-obs
%   CHI: colonnes-source-doc, colonnes-source-hyp, colonnes-source-obs, colonnes-source-registre, comparaisons-modele, creances-part-au-dela-un-an
%   CHI: doublons-part-mesuree, facturation-part-organisme, facturation-part-patient, feries-urgences-hors-moyenne, feries-urgences-moyenne, flux-lundi-moyenne
%   CHI: flux-programme-moyenne, flux-samedi-moyenne, hebdo-rapport-mesure, injecte-correlation-parametres, instantane-duree-rafraichissement, instantane-erreurs-completude
%   CHI: instantane-erreurs-dimension, instantane-fenetre-completude, instantane-fenetre-dimension, instantane-lectures-temoin, instantane-lignes, instantane-objets-copies
%   CHI: instantane-rapport-cout-lecture, instantane-rapport-hors-completude, instantane-tables, instantane-volume, jour-meme-part-marge-nulle, niveaux-tri
%   CHI: paires-candidates, paires-possibles, part-couverture-principale, part-jour-meme, part-sans-couverture, part-sejours-depuis-urgences
%   CHI: part-sexe-masculin, part-tri-niveau-4, part-tri-niveau-5, part-urbain, qualite-taux-quarantaine-maximal, ramadan-baisse
%   CHI: ramadan-hors-moyenne, ramadan-programme-moyenne, rapport-accord-etat, rapport-accord-etat-naissance, rapport-accord-pays-naissance, recouvrement-taux-aboutissement
%   CHI: recouvrement-taux-global, recouvrement-taux-maximal, recouvrement-taux-minimal, relation-brute-complement, relation-brute-ensemble, relation-inter-complement
%   CHI: relation-inter-ensemble, relation-intra-complement, relation-intra-ensemble, source-part-non-appropriees, taches-graphe, taux-absenteisme-global
%   CHI: taux-annulation-global, taux-encaissement, tdb-erreurs-lecture-tolerees, tranches-age, urgences-part-consultation-ordinaire, urgences-part-hospitalisation
%   CHI: urgences-part-mode-principal, urgences-part-retour-domicile, urgences-part-transfert, urgences-rapport-pointe-creux
%   SER: facturation-par-type, instantane-objets, qualite-par-table
%   SECTIONS: 9

\chapitreDAnnexe{Compléments techniques}\label{annexe:complements}

\section{La provenance des colonnes}\label{sec:complements-provenance}

Chaque colonne de la couche source est rattachée à ce qui l'établit. Le
tableau~\ref{tab:correspondance-champs} donne sa correspondance avec les écrans pour deux tables. Le
dictionnaire complet la porte pour toutes les tables, et l'annexe~\ref{annexe:dictionnaire} en donne
la synthèse. La colonne « Preuve » porte l'identifiant du relevé, celui que le registre des champs
référence.

% SÉCABLE : en `tabular`, le tableau partait d'un seul bloc à la page suivante dès qu'il ne tenait
% pas sur la fin de la page courante, et laissait une page presque blanche. En `longtable`, il
% commence où il se trouve et répète sa ligne d'intitulés en tête de chaque page qu'il occupe.
%
% Les largeurs sont MESURÉES sur le rendu, non estimées : `type_piece_identite` occupe 4,18 cm en
% chasse fixe et surimprimait la colonne suivante, et l'intitulé « Provenance » en occupe 2,06. La
% somme des cinq colonnes vaut 12,9 cm.
\begin{longtable}{@{}p{2.6cm}p{4.3cm}p{1.0cm}p{2.1cm}p{2.9cm}@{}}
\caption[Correspondance entre libellé d'écran, colonne et preuve d'observation]{Correspondance entre libellé d'écran, colonne et preuve d'observation, pour deux tables}\label{tab:correspondance-champs}\\
\toprule
Libellé relevé & Colonne & Type & Provenance & Preuve \\
\midrule
\endfirsthead
\toprule
Libellé relevé & Colonne & Type & Provenance & Preuve \\
\midrule
\endhead
\multicolumn{5}{@{}l}{\textbf{\texttt{patients}} — extraits}\\
N° IPP & \texttt{n\_ipp} & texte & \texttt{OBS} & \texttt{REL-PAT.D01} \\
Nom & \texttt{nom} & texte & \texttt{OBS} & \texttt{REL-PAT.D02} \\
Nom de famille 1 & \texttt{nom\_famille\_1} & texte & \texttt{OBS} & \texttt{REL-PAT.D03} \\
Sexe & \texttt{sexe} & texte & \texttt{OBS} & \texttt{REL-PAT.D05} \\
D. Nai & \texttt{date\_naissance} & texte & \texttt{OBS} & \texttt{REL-PAT.D06} \\
Type pièce d'identité & \texttt{type\_piece\_identite} & texte & \texttt{OBS} & \texttt{REL-PAT.D09} \\
N° pièce d'identité & \texttt{n\_piece\_identite} & texte & \texttt{OBS} & \texttt{REL-PAT.D10} \\
E. Civil & \texttt{etat\_civil} & texte & \texttt{OBS} & \texttt{REL-PAT.D11} \\
Compagnie d'assur. & \texttt{compagnie\_assurance} & texte & \texttt{OBS} & \texttt{REL-PAT.A01} \\
N° Assu & \texttt{n\_assure} & texte & \texttt{OBS} & \texttt{REL-PAT.A03} \\
Ville & \texttt{ville} & texte & \texttt{OBS} & \texttt{REL-PAT.H05} \\
Nationalité & \texttt{nationalite} & texte & \texttt{OBS} & \texttt{REL-PAT.H07} \\
Avertissements SMS & \texttt{avertissements\_sms} & texte & \texttt{OBS} & \texttt{REL-PAT.H12} \\
\addlinespace
\multicolumn{5}{@{}l}{\textbf{\texttt{rendez\_vous}} — extraits}\\
Agenda & \texttt{agenda} & texte & \texttt{OBS} & \texttt{REL-RDV.R01} \\
Activité & \texttt{activite} & texte & \texttt{OBS} & \texttt{REL-RDV.R02} \\
Origine & \texttt{origine} & texte & \texttt{OBS} & \texttt{REL-RDV.R03} \\
Date rendez-vous & \texttt{date\_heure\_rdv} & texte & \texttt{OBS} & \texttt{REL-RDV.R08} \\
Type d'attention & \texttt{type\_attention} & texte & \texttt{OBS} & \texttt{REL-RDV.R10} \\
État & \texttt{etat} & texte & \texttt{OBS} & \texttt{REL-RDV.R11} \\
Durée & \texttt{duree} & texte & \texttt{OBS} & \texttt{REL-RDV.R12} \\
Créé par & \texttt{cree\_par} & texte & \texttt{OBS} & \texttt{REL-RDV.C01} \\
Date création & \texttt{date\_creation} & texte & \texttt{OBS} & \texttt{REL-RDV.C02} \\
Confirmé par & \texttt{confirme\_par} & texte & \texttt{OBS} & \texttt{REL-RDV.C05} \\
Annulé par & \texttt{annule\_par} & texte & \texttt{OBS} & \texttt{REL-RDV.C07} \\
Date annul. & \texttt{date\_annulation} & texte & \texttt{OBS} & \texttt{REL-RDV.C08} \\
\bottomrule
\end{longtable}

Le registre des champs porte \chiffre{colonnes-source-registre} entrées, et le catalogue de la base
compte exactement \chiffre{colonnes-source-catalogue} colonnes au schéma source. Un contrôle bloquant
vérifie la correspondance dans les deux sens, et le décompte par étiquette se lit des deux côtés
(voir tableau~\ref{tab:provenance-colonnes}). Les couches intermédiaire et analytique comptent ensemble \chiffre{colonnes-aval-total} colonnes,
\chiffre{colonnes-intermediate} et \chiffre{colonnes-marts}, et aucune ne porte d'étiquette de
provenance. Une colonne calculée a pour provenance sa formule, écrite dans le modèle qui la produit.
Ces \chiffre{colonnes-aval-decrites} colonnes sont toutes décrites dans les fichiers de propriétés de
l'outil de transformation.

\begin{table}[htbp]
\centering
\caption{Provenance des colonnes de la couche source, comptée au registre et au catalogue de la base}\label{tab:provenance-colonnes}
\begin{tabular}{@{}p{2.2cm}rrp{6.4cm}@{}}
\toprule
Étiquette & Registre & Catalogue & Ce qu'elle signifie \\
\midrule
\texttt{OBS} & \chiffre{colonnes-source-obs} & \chiffre{colonnes-source-catalogue-obs} & le champ a été vu à l'écran, et son libellé est reproduit \\
\texttt{DOC} & \chiffre{colonnes-source-doc} & \chiffre{colonnes-source-catalogue-doc} & le champ est déduit d'une source documentaire \\
\texttt{HYP} & \chiffre{colonnes-source-hyp} & \chiffre{colonnes-source-catalogue-hyp} & le champ est posé, et l'absence de preuve est écrite \\
\midrule
\textbf{Total} & \textbf{\chiffre{colonnes-source-registre}} & \textbf{\chiffre{colonnes-source-catalogue}} & \\
\bottomrule
\end{tabular}
\end{table}

\section{Les distributions du générateur}\label{sec:complements-distributions}

Le tableau~\ref{tab:distributions} donne chaque distribution du générateur, sa valeur ou sa forme,
son étiquette et ce qui l'étaye. La structure par âge et la répartition par sexe sont des
conventions posées : une conclusion tirée de la pyramide des âges ne dirait rien de la population
réelle.

% SÉCABLE POUR LA MÊME RAISON QUE LE TABLEAU DE CORRESPONDANCE : rendu insécable, il partait d'un
% bloc et laissait une demi-page blanche.
\begin{longtable}{@{}p{4.6cm}p{4.4cm}p{1.7cm}p{2.6cm}@{}}
\caption[Les distributions du générateur]{Les distributions du générateur et l'étiquette de provenance de chacune}\label{tab:distributions}\\
\toprule
Distribution & Valeur ou forme & Étiquette & Ce qui l'étaye \\
\midrule
\endfirsthead
\toprule
Distribution & Valeur ou forme & Étiquette & Ce qui l'étaye \\
\midrule
\endhead
Structure par âge & \chiffre{tranches-age} tranches & \texttt{HYP} & convention posée \\
Répartition par sexe & \chiffre{part-sexe-masculin} d'hommes & \texttt{HYP} & convention posée \\
Environnement de résidence & \chiffre{part-urbain} en milieu urbain & \texttt{HYP} & mesure régionale \cite{s30} \\
Couverture principale & \chiffre{part-couverture-principale} & \texttt{HYP} & répartition nationale des régimes \\
Absence de couverture & \chiffre{part-sans-couverture} & \texttt{HYP} & convention posée \\
Niveaux de tri aux urgences & \chiffre{niveaux-tri} niveaux & \texttt{DOC} & échelle de tri publiée \\
Part du niveau de tri 4 & \chiffre{part-tri-niveau-4} & \texttt{DOC} & échelle de tri publiée \\
Part du niveau de tri 5 & \chiffre{part-tri-niveau-5} & \texttt{DOC} & échelle de tri publiée \\
Séjours provenant des urgences & \chiffre{part-sejours-depuis-urgences} & \texttt{HYP} & convention posée \\
\bottomrule
\end{longtable}

\section{L'absence d'activité chirurgicale}\label{sec:complements-chirurgie}

Trois vérifications fondent la lecture du recueil statistique \cite{s30}. Le nom de
l'établissement, cherché sans égard à la casse dans le texte intégral des deux éditions, ne figure
nulle part dans l'emprise du tableau des interventions. La somme des lignes visibles de la section
provinciale égale exactement le total imprimé, ce qui exclut un établissement non listé. Ces
tableaux recensent enfin une activité, et non un parc d'établissements : y être absent signifie ne
pas l'exercer.

\section{Les outils de la chaîne}\label{sec:complements-pile}

La figure~\ref{fig:pile-technique} donne chaque brique de la pile, son rôle et la couche à laquelle
elle appartient.

\begin{figure}[htbp]
\centering
\begin{tikzpicture}[x=1mm, y=1mm]

\node[coucheUML, anchor=east] at (-4,0)   {restitution};
\node[coucheUML, anchor=east] at (-4,-20) {ordonnancement};
\node[coucheUML, anchor=east] at (-4,-40) {identités};
\node[coucheUML, anchor=east] at (-4,-60) {modélisation};
\node[coucheUML, anchor=east] at (-4,-80) {persistance};

\node[briqueUML, text width=52mm] (res) at (26,0)   {Bibliothèque d'application web analytique\\[1pt]\emph{lire, filtrer, afficher}};
\node[briqueUML, text width=52mm] (ord) at (26,-20) {Orchestrateur de graphes de tâches\\[1pt]\emph{déclencher dans l'ordre}};
\node[briqueUML, text width=52mm] (rap) at (26,-40) {Bibliothèque de rapprochement probabiliste\\[1pt]\emph{estimer, apparier, regrouper}};
\node[briqueUML, text width=52mm] (mod) at (26,-60) {Outil de transformation déclaratif\\[1pt]\emph{typer, modéliser, tester}};
\node[briqueUML, text width=52mm] (per) at (26,-80) {Serveur de base de données relationnelle\\[1pt]\emph{stocker, servir, transiger}};

\node[briqueOptUML, text width=44mm] (sch) at (96,-60) {Schéma en étoile\\[1pt]\emph{faits, dimensions, agrégats}};
\node[briqueOptUML, text width=44mm] (mem) at (96,-40) {Moteur en mémoire\\[1pt]\emph{calculer hors du serveur}};

\draw[fluxUML] (per) -- (mod);
\draw[fluxUML] (mod) -- (rap);
\draw[fluxUML] (rap) -- (ord);
\draw[fluxUML] (ord) -- (res);
\draw[fluxUML] (sch.west) -- (mod.east);
\draw[fluxUML] (mem.west) -- (rap.east);
\end{tikzpicture}
\caption{La pile technique, couche par couche}\label{fig:pile-technique}
\end{figure}

Chaque brique porte sa désignation fonctionnelle et son rôle en trois verbes, et le
chapitre~\ref{chap:outils} nomme l'outil qui l'incarne. Les flèches pleines montent d'une couche à
la suivante. Les deux briques en pointillé sont des choix internes à la couche qu'elles rejoignent.

\section{La couche analytique et l'orchestration}\label{sec:complements-analytique}

La figure~\ref{fig:schema-etoile} donne les dimensions et les faits de la couche analytique,
sous les noms que le code leur donne. Les agrégats, non représentés, se calculent depuis les
faits.

\begin{figure}[htbp]
\centering
\begin{tabular}{@{}c@{}}
\fbox{\begin{minipage}{0.92\linewidth}\centering\small
\fbox{\begin{minipage}{0.28\linewidth}\centering \texttt{dim\_date}\end{minipage}}%
\hfill
\fbox{\begin{minipage}{0.28\linewidth}\centering \texttt{dim\_patient} \\ \footnotesize historisée\end{minipage}}%
\hfill
\fbox{\begin{minipage}{0.28\linewidth}\centering \texttt{dim\_service}\end{minipage}}\\[3pt]
% CHAQUE FLÈCHE EST CENTRÉE DANS UNE BOÎTE DE LA LARGEUR DE LA DIMENSION QU'ELLE RELIE. Avec un
% `\hfill` entre elles, les deux flèches obliques étaient rejetées aux bords du cadre et
% s'arrêtaient dans le vide, à côté des boîtes qu'elles devaient joindre. Les trois dimensions
% occupent chacune 0,28 de la largeur ; les trois boîtes de flèche prennent la même largeur.
\makebox[0.28\linewidth]{$\searrow$}\hfill\makebox[0.28\linewidth]{$\downarrow$}\hfill\makebox[0.28\linewidth]{$\swarrow$}\\[3pt]
\fbox{\begin{minipage}{0.94\linewidth}\centering \textbf{Les six faits} \\[2pt]
\footnotesize
\texttt{fct\_sejour} — un séjour \quad
\texttt{fct\_passage} — un passage, les trois types \\
\texttt{fct\_passage\_urgence} — un passage aux urgences \quad
\texttt{fct\_rendez\_vous} — un rendez-vous \\
\texttt{fct\_facturation} — une facture \quad
\texttt{fct\_encaissement} — un encaissement
\end{minipage}}\\[3pt]
\makebox[0.28\linewidth]{$\nearrow$}\hfill\makebox[0.28\linewidth]{$\uparrow$}\hfill\makebox[0.28\linewidth]{$\nwarrow$}\\[3pt]
\fbox{\begin{minipage}{0.28\linewidth}\centering \texttt{dim\_activite}\end{minipage}}%
\hfill
\fbox{\begin{minipage}{0.28\linewidth}\centering \texttt{dim\_agent}\end{minipage}}%
\hfill
\fbox{\begin{minipage}{0.28\linewidth}\centering \texttt{dim\_organisme}\end{minipage}}
\end{minipage}}
\end{tabular}
\caption{Le schéma en étoile de la couche analytique}\label{fig:schema-etoile}
\end{figure}

L'un des agrégats est une vue posée hors de l'outil de transformation, par un fichier de définition
de schéma appliqué avant l'initialisation. Il n'est donc pas l'un des modèles, et son fichier de
propriétés donne le motif de ce maintien à part. L'outil de transformation le signale d'ailleurs à
l'analyse, faute de modèle qui corresponde à cet objet documenté.

Le grain de chaque fait est déclaré. Le fait des séjours porte un séjour par ligne et agrège la
couche des mouvements, qui porte une admission, au plus un transfert et au plus une sortie. Le fait
des passages couvre les trois types d'épisode, et c'est à lui que se raccordent les factures et les
prises en charge. Le fait des passages aux urgences en est un satellite, dont toutes les clés sont
des clés du fait des passages. Le grain du fait des encaissements est l'encaissement et non la
facture, une facture pouvant recevoir plusieurs règlements.

La clé d'une version de la dimension patient est le couple formé par l'identifiant et le début de
validité. Les intervalles sont semi-ouverts : le début est inclus, la fin est exclue, et la version
courante n'a pas de fin. Deux versions consécutives se touchent donc exactement, sans chevauchement
ni trou. Un fait daté pointe vers ce couple, donc vers la version en vigueur au moment de
l'événement.

Trois contrôles vérifient ces versions à chaque construction. La fin d'une version non courante
égale le début de la suivante, deux versions ne se chevauchent pas, et un identifiant porte une
seule version courante. La continuité seule laisserait passer deux versions courantes, et le
non-chevauchement seul laisserait passer un trou.

{\raggedright Les \chiffre{taches-graphe} tâches du graphe portent au code les identifiants
suivants, dans l'ordre de la figure~\ref{fig:sequence-quotidienne} :
\texttt{verifier\_disponibilite\_extraction}, \texttt{charger\_journee}, \texttt{controle\_qualite},
\texttt{dbt\_intermediaire}, \texttt{dbt\_dimensions}, \texttt{dbt\_faits}, \texttt{dbt\_tests},
\texttt{rapprochement\_prediction}, \texttt{rapprochement\_regroupement\_evaluation},
\texttt{dbt\_agregats}, \texttt{rafraichir\_instantane}, \texttt{exporter}.\par}

\section{Le modèle de rapprochement}\label{sec:complements-modele}

Deux probabilités donnent son poids à chaque champ, et leur rapport fait le poids selon la formule
de la littérature \cite{s22}. La première est celle d'observer un accord quand les deux fiches
désignent la même personne. Elle mesure la fiabilité du champ : un téléphone change, une date de
naissance est stable. La seconde est celle d'observer un accord entre deux personnes différentes,
et elle mesure le hasard : deux personnes partagent souvent une ville, rarement une date de
naissance.

Un désaccord porte aussi de l'information. Sur un champ stable, il retire beaucoup de poids, et sur
un champ instable il en retire peu. L'addition des poids suppose les champs indépendants une fois
l'état de la paire connu.

Ils ne le sont pas, car une adresse et une ville sont liées, comme un nom
de père et un nom de famille. Les poids extrêmes sont donc surestimés, et les probabilités
affichées plus tranchées que la réalité. Le classement des paires n'en souffre pas, et c'est ce dont
le seuil a besoin, mais une probabilité isolée ne se lit pas comme une chance.

Le tableau~\ref{tab:regles-blocage} donne les quatre règles de blocage et les paires que chacune
produit. La réduction y est rapportée aux \chiffre{paires-possibles} paires possibles.

\begin{table}[htbp]
\centering
\caption{Les règles de blocage, les paires que chacune produit et la réduction qu'elle opère}\label{tab:regles-blocage}
\begin{tabular}{@{}p{6.6cm}rr@{}}
\toprule
Règle de blocage & Paires & Réduction \\
\midrule
premier nom de famille \emph{et} adresse & \chiffre{blocage-nom-adresse} & \chiffre{blocage-reduction-nom-adresse}~\% \\
premier nom de famille \emph{et} téléphone & \chiffre{blocage-nom-telephone} & \chiffre{blocage-reduction-nom-telephone}~\% \\
type \emph{et} numéro de pièce d'identité & \chiffre{blocage-piece-identite} & \chiffre{blocage-reduction-piece-identite}~\% \\
nom du père \emph{et} nom de la mère \emph{et} date de naissance & \chiffre{blocage-parents-naissance} & \chiffre{blocage-reduction-parents-naissance}~\% \\
\midrule
\textbf{Union des quatre} & \textbf{\chiffre{paires-candidates}} & \textbf{\chiffre{blocage-reduction-union}~\%} \\
\bottomrule
\end{tabular}
\end{table}

Treize colonnes sont comparées sous \chiffre{comparaisons-modele} comparaisons, car les deux
colonnes de nom de famille se traitent ensemble. Les champs de texte utilisent une distance de
Jaro-Winkler à deux seuils successifs. La date de naissance passe par l'égalité exacte, puis par un
écart en nombre de jours, puis par une distance de chaîne, puis par le niveau par défaut. L'ordre
des deux niveaux intermédiaires a été fixé par la mesure : un écart en jours sépare nettement mieux
les paires qu'une distance de chaîne.

Neuf colonnes sont écartées, chacune pour un motif mesuré. Pour certaines, un accord n'apprend
presque rien, car le rapport entre la probabilité d'accord chez les vrais couples et chez les faux vaut
\chiffre{rapport-accord-etat} pour l'état, \chiffre{rapport-accord-pays-naissance} pour le pays de
naissance et \chiffre{rapport-accord-etat-naissance} pour l'état de naissance. D'autres colonnes sont des
copies d'une autre colonne faites par le programme, égales ligne à ligne sur toute la population.
L'état de naissance copie ainsi l'état, et la ville de naissance copie la ville. Les comparer compterait
deux fois le même accord.

\section{Deux tableaux de l'analyse de l'activité}\label{sec:complements-analyse}

La facturation et la quarantaine se détaillent en deux tableaux. Le premier ventile le montant
facturé par type d'épisode (cf. tableau~\ref{tab:montants-episode}).

\begin{table}[htbp]
\centering
\caption{Montants facturés par type d'épisode, en dirhams}\label{tab:montants-episode}
\footnotesize
\pgfplotstabletypeset[
  col sep=comma,
  columns/episode/.style={column name={Type d'épisode}, string type},
  columns/factures/.style={column name={Factures}, int detect, 1000 sep={\,}},
  columns/montant/.style={column name={Montant}, fixed, fixed zerofill, precision=2,
    1000 sep={\,}},
  columns/part/.style={column name={Part (\%)}, fixed, fixed zerofill, precision=2},
  columns/organisme/.style={column name={Part organisme (\%)}, fixed, fixed zerofill,
    precision=2},
  every head row/.style={before row=\toprule, after row=\midrule},
  every last row/.style={after row=\bottomrule},
]{\serie{facturation-par-type}}
\end{table}

Les trois codes d'épisode s'y affichent sans libellé : la taxonomie réglementaire des modes
d'utilisation de l'hôpital est établie sans être énumérée \cite{s27}.

Le second donne la complétude et la quarantaine de chaque table de la couche intermédiaire (cf.
tableau~\ref{tab:qualite-par-table}).

\begin{table}[htbp]
\centering
\caption{Complétude et quarantaine par table de la couche intermédiaire}\label{tab:qualite-par-table}
\footnotesize
% L'espacement entre colonnes est resserré POUR CE TABLEAU SEUL, et la mesure dit de combien : à
% 6 pt, l'intitulé abrégé laissait encore 0,33 cm de débordement ; à 4 pt, les six colonnes
% tiennent dans le bloc de texte, que le filet d'en-tête de page donne exactement. Le groupe
% `center` borne la portée du changement.
\setlength{\tabcolsep}{4pt}
\pgfplotstabletypeset[
  col sep=comma,
  % Les noms de table portent des soulignements, qui sont des caractères actifs en mode texte.
  % `string replace*` les échappe à la lecture, sans toucher au fichier de données.
  %
  % LE NOM DE TABLE SE COMPOSE EN CHASSE FIXE, comme partout ailleurs dans le document. Il était
  % composé en romain, où `\textunderscore` produit une barre basse démesurée : `int_patients` s'y
  % lisait `int__patients`. Le défaut ne se voyait qu'à la relecture du document composé.
  columns/objet/.style={column name={Table}, string type, string replace*={_}{\_},
    postproc cell content/.append style={@cell content/.add={\texttt\bgroup}{\egroup}}},
  columns/colonnes/.style={column name={Colonnes}, int detect},
  columns/incompletes/.style={column name={Incomplètes}, int detect},
  % L'INTITULÉ EST ABRÉGÉ PARCE QUE LE TABLEAU DÉBORDAIT LA MARGE, mesuré au rendu : son filet
  % s'étendait 1,05 cm au-delà du bloc de texte, que le filet d'en-tête de page donne exactement.
  % La légende sous le tableau porte le mot entier.
  columns/completude_minimale/.style={column name={Complétude min. (\%)}, fixed,
    fixed zerofill, precision=2},
  columns/quarantaine/.style={column name={Quarantaine}, int detect, 1000 sep={\,}},
  columns/taux_quarantaine/.style={column name={Taux (\%)}, fixed, fixed zerofill, precision=2},
  every head row/.style={before row=\toprule, after row=\midrule},
  every last row/.style={after row=\bottomrule},
]{\serie{qualite-par-table}}
\end{table}

La colonne « Complétude min. » y porte la complétude minimale observée sur les colonnes de la table.

\section{Les valeurs de l'analyse en précision complète}\label{sec:complements-precision}

L'analyse de l'activité affiche ses pourcentages et ses moyennes à une décimale, et ses coefficients
à deux décimales. Chaque valeur affichée a sa propre entrée au registre des chiffres, dont la
commande arrondit celle de la valeur complète. Le tableau~\ref{tab:precision-complete} donne ces
valeurs en précision complète, dans l'ordre où l'analyse les emploie.

{\footnotesize
\begin{longtable}{@{}p{0.72\linewidth}r@{}}
\caption[Les valeurs de l'analyse en précision complète]{Les valeurs de l'analyse de l'activité en précision complète}\label{tab:precision-complete}\\
\toprule
Grandeur & Valeur \\
\midrule
\endfirsthead
\toprule
Grandeur & Valeur \\
\midrule
\endhead
\bottomrule
\endlastfoot
Passages programmés par jour ouvert & \chiffre{flux-programme-moyenne} \\
Passages programmés par jour, premier jour de la semaine & \chiffre{flux-lundi-moyenne} \\
Passages programmés par jour, sixième jour de la semaine & \chiffre{flux-samedi-moyenne} \\
Rapport entre ces deux moyennes & \chiffre{hebdo-rapport-mesure} \\
Passages programmés par jour ouvert pendant le Ramadan & \chiffre{ramadan-programme-moyenne} \\
Passages programmés par jour ouvert hors Ramadan & \chiffre{ramadan-hors-moyenne} \\
Baisse pendant le Ramadan & \chiffre{ramadan-baisse}~\% \\
Passages aux urgences par jour férié & \chiffre{feries-urgences-moyenne} \\
Passages aux urgences par jour non férié & \chiffre{feries-urgences-hors-moyenne} \\
Délais égaux à la marge disponible & \chiffre{artefact-part}~\% \\
Rendez-vous pris pour le jour même & \chiffre{part-jour-meme}~\% \\
Part de ces rendez-vous à marge nulle ou négative & \chiffre{jour-meme-part-marge-nulle}~\% \\
Taux d'absentéisme & \chiffre{taux-absenteisme-global}~\% \\
Taux d'annulation & \chiffre{taux-annulation-global}~\% \\
Corrélation entre activités, ensemble & \chiffre{relation-inter-ensemble} \\
Corrélation à l'intérieur des activités, ensemble & \chiffre{relation-intra-ensemble} \\
Corrélation sans distinguer les activités, ensemble & \chiffre{relation-brute-ensemble} \\
Corrélation entre activités, complément & \chiffre{relation-inter-complement} \\
Corrélation à l'intérieur des activités, complément & \chiffre{relation-intra-complement} \\
Corrélation sans distinguer les activités, complément & \chiffre{relation-brute-complement} \\
Corrélation entre les deux tables de paramètres & \chiffre{injecte-correlation-parametres} \\
Rapport entre pointe et creux horaires des urgences & \chiffre{urgences-rapport-pointe-creux} \\
Retour à domicile après un passage aux urgences & \chiffre{urgences-part-retour-domicile}~\% \\
Transfert & \chiffre{urgences-part-transfert}~\% \\
Hospitalisation & \chiffre{urgences-part-hospitalisation}~\% \\
Mode d'arrivée le plus fréquent & \chiffre{urgences-part-mode-principal}~\% \\
Deux niveaux de tri les moins graves & \chiffre{urgences-part-consultation-ordinaire}~\% \\
Part prise en charge par un organisme & \chiffre{facturation-part-organisme}~\% \\
Part restant au patient & \chiffre{facturation-part-patient}~\% \\
Taux de recouvrement global & \chiffre{recouvrement-taux-global}~\% \\
Taux de recouvrement le plus bas, par débiteur & \chiffre{recouvrement-taux-minimal}~\% \\
Taux de recouvrement le plus haut, par débiteur & \chiffre{recouvrement-taux-maximal}~\% \\
Taux d'encaissement & \chiffre{taux-encaissement}~\% \\
Restant dû porté par des créances d'un an et plus & \chiffre{creances-part-au-dela-un-an}~\% \\
Relances abouties & \chiffre{recouvrement-taux-aboutissement}~\% \\
Taux de quarantaine le plus élevé d'une table & \chiffre{qualite-taux-quarantaine-maximal}~\% \\
Personnes portant deux fiches & \chiffre{doublons-part-mesuree}~\% \\
Consultations non appropriées dans l'étude provinciale, en proportion & \chiffre{source-part-non-appropriees} \\
\end{longtable}
}

\section{Le motif de l'instantané du tableau de bord}\label{sec:complements-instantane}

Le tableau de bord lit un schéma de tables figées parce qu'un lecteur voyait une erreur pendant la
reconstruction, et non pour le coût des lectures. Ce coût a été mesuré : lire les indicateurs de
toutes les pages contre les vues coûte \chiffre{instantane-rapport-cout-lecture} fois ce qu'il coûte
contre des tables copiées. Ce rapport moyen est porté par une minorité d'indicateurs. Sans la page
la plus coûteuse, il tombe à \chiffre{instantane-rapport-hors-completude}, un gain trop inégal pour
justifier un schéma de plus à maintenir.

Un lecteur concurrent a lu en boucle pendant une reconstruction complète. Il a compté
\chiffre{instantane-erreurs-dimension} lectures en erreur sur \chiffre{instantane-lectures-temoin}
tentatives contre une dimension simple, et \chiffre{instantane-erreurs-completude} contre la vue la
plus coûteuse à produire. Le message disait que l'objet n'existait plus, pendant
\chiffre{instantane-fenetre-dimension} seconde dans le premier cas et
\chiffre{instantane-fenetre-completude} dans le second.

Le critère retenu traite différemment les deux échecs. Le nombre d'échecs de lecture tolérés est de
\chiffre{tdb-erreurs-lecture-tolerees}, sous toute charge, parce qu'un utilisateur les subit sans
recours. Un échec de remplacement est journalisé et relançable. La transaction qui l'enveloppe
laisse alors l'instantané dans sa génération précédente, complète.

Les objets que les pages lisent se rangent par effectif (voir figure~\ref{fig:objets-instantane}).

\begin{figure}[htbp]
\centering
\begin{tikzpicture}
\begin{axis}[
  width=11cm, height=4.6cm, ymode=log,
  xlabel={objets copiés, rangés par effectif décroissant},
  ylabel={lignes copiées},
  xmin=0, xmax=25,
  grid=major, grid style={gray!25},
  every axis plot/.append style={thick},
]
\addplot[color=black, mark=*, mark size=1.2pt] table[x=rang, y=lignes, col sep=comma]
  {\serie{instantane-objets}};
\end{axis}
\end{tikzpicture}
\caption{Les objets que le tableau de bord lit, par effectif}\label{fig:objets-instantane}
\end{figure}

La figure range les \chiffre{instantane-objets-copies} objets par effectif, et son abscisse est un
rang. L'instantané tient en \chiffre{instantane-tables} tables : ces copies, et deux tables de
service qui portent l'état et les paramètres. Il pèse \chiffre{instantane-volume} octets pour
\chiffre{instantane-lignes} lignes, et se reconstruit en \chiffre{instantane-duree-rafraichissement}
secondes. Les effectifs s'étalent sur quatre ordres de grandeur, du fait de passage à l'agrégat de
doublons. Cet écart explique que le gain de lecture varie autant d'un indicateur à l'autre.

Le registre des indicateurs déclare des objets de la couche analytique, quand les pages lisent le
schéma figé. Le registre nomme la couche où un modèle produit l'objet, et le nom de table, privé de
son schéma, relie les deux.
